\documentclass[11pt]{article}
\usepackage[T1]{fontenc}
\usepackage{lmodern,amsmath,amssymb,amsthm,mathtools,booktabs,geometry,microtype,xcolor,enumitem,hyperref,xurl}
\geometry{left=0.94in,right=0.94in,top=0.88in,bottom=0.88in}
\definecolor{ink}{HTML}{18344D}
\hypersetup{colorlinks=true,linkcolor=ink,urlcolor=ink,citecolor=ink,pdftitle={Positive Frobenius deficits and a finite-difference route beyond Dirichlet GRH}}
\newtheorem{theorem}{Theorem}[section]
\newtheorem{lemma}[theorem]{Lemma}
\newtheorem{corollary}[theorem]{Corollary}
\newtheorem{proposition}[theorem]{Proposition}
\theoremstyle{definition}\newtheorem{definition}[theorem]{Definition}
\newcommand{\Q}{\mathbb Q}\newcommand{\C}{\mathbb C}\newcommand{\R}{\mathbb R}\newcommand{\Rea}{\operatorname{Re}}\newcommand{\Frob}{\operatorname{Frob}}\newcommand{\Gal}{\operatorname{Gal}}\newcommand{\pos}[1]{\left(#1\right)_+}
\title{\color{ink}\bfseries Positive Frobenius deficits and a finite-difference route beyond Dirichlet GRH\\[6pt]\large Transfer analysis of the Alp\"oge--Claude--Sernau proof}
\author{Research memorandum for ProveIt}
\date{25 November 2026}
\begin{document}
\maketitle
\begin{abstract}
The compact Lean proof of RH and Dirichlet GRH combines a comparison of covolumes of theta lifts with a one-sided Mellin--Landau detector. We prove an extension of the analytic detector using a nonnegative Frobenius deficit, a fixed multiplicative finite difference, and a known abelian baseline. The difference removes an unknown centering constant, reverses the real-pole sign correctly, and allows an averaged positive-part bound in place of pointwise control. An optimal nonnegative lift of the deficit modulo abelian class functions is given for arbitrary permutation representations. For $S_4$, two Frobenius classes suffice and are necessary among detectors with strictly negative nonabelian Fourier coefficients. Two low-degree subfields realize the detector without requiring unknown Artin holomorphy. For the Kummer fields in Hooley's theorem, a single unweighted congruence-and-power-residue test suffices per field. The same detector applies to normalized holomorphic modular forms and their symmetric powers. These are proved reductions to explicit inequalities, not proofs of the new Riemann hypotheses. The missing arithmetic-geometric inequality is isolated; no priority claim is made.
\end{abstract}
\tableofcontents
\newpage
\section{Scope and provenance}
The repository \cite{repo} reports a compact edition (about 76,000 lines of code) and a readable edition (about 300,000 lines), both with the same five conclusions. The inspected compact tree is identified by Git tree object \nolinkurl{72fead102f663d5c7a407964a13dc0cf9d977207}. We pin individual file blobs in \texttt{source-audit.tsv}. The tree identifier is what the GitHub tree endpoint returned; it need not be the identifier of a commit object. The source uses Lean 4.35.0 and Mathlib pinned to \nolinkurl{56a005855ea17e0b0958f73ec801ae140b0a1679}.

The five statements, isolated in \texttt{src/Statements.lean} with only Mathlib imported, cover the Riemann zeta function and Dirichlet characters of every modulus. The right-half-plane statement excludes the pole of the principal character at $s=1$. They do not quantify over arbitrary Dedekind zeta-functions, Artin representations or automorphic forms. The Lean FRO and Mathlib Initiative certified the original proof and statement of RH on November 21 \cite{cert}; Rubinstein reports independent checking of the original and the shortened proof \cite{rubinstein}. We use Dirichlet GRH as an established input, without claiming that the human exposition is complete or that the present environment reproduced the Lean builds. The repository's README reports kernel and comparator checks. No local Lean executable was available.

\paragraph{Earlier ProveIt reports.} The incoming directory \cite{prior} was inspected with the GitHub connector; four ZIP files were retrieved in base64 and selected README and TeX passages extracted in memory. In particular, \nolinkurl{grh_relative_prime_races_2026-11-25.zip} already establishes signed, constant-free relative-zeta criteria, normal tower compression and a mixed compositum criterion. Its $S_4$ example uses weights $8,-1$ on completely split and $1\,3$ primes. \nolinkurl{grh_transfers.zip} gives a cubic Artin sign test; \nolinkurl{grh_method_transfers_2026-11-25.zip} contains a one-sided Mellin detector and a transfer obstruction; \nolinkurl{grh-arithmetic-transfers.zip} reports further arithmetic results and a four-weight Dedekind criterion. We have not independently audited those results. Our particular additions are a nonnegative optimal representative modulo abelian class functions, an averaged one-sided finite-difference detector, the two-class $S_4$ construction using degree four and six fields, and its Kummer and modular applications. Other incoming archives were not exhaustively inspected; overlap or earlier priority is possible.

\section{The source mechanism and its reusable interface}
Write $R_{d,M}(x)$ for the source's weighted reciprocal ramp. Its arithmetic coefficient is
\[
 \lambda(n)=\Lambda(n)\,\mathbf 1_{n\equiv1\pmod M}(1-\chi_d(n)),
\]
where $\chi_d$ is the quadratic character attached to the chosen imaginary quadratic field. The function \texttt{primeError} subtracts $\log x/\varphi(M)$ and a logarithmic-derivative constant from $\sum_{n\le x}\lambda(n)(1/n-1/x)$. It is rewritten as an exact finite divisibility sum in \texttt{New/Endgame/PrimeError.lean} (lines 25--69).

The central comparison is \texttt{star\_inequality} in \texttt{New/Main.lean}, line 121, which gives
\[
 m\,\mathrm{remainder}\le (\log b)^2\,\mathrm{budget},\qquad
 \mathrm{budget}=fa(\log b)^3.
\]
Here $m$ is the number of lifts, of order $\sqrt x$, and $f,a$ count places and scale the arithmetic order. The proof compares the same squared covolume $\exp V$ and lattice index bounded by $\exp\delta$ on a nearly spanning family. With $d_0$ the row dimension, the schematic inequality is
\[
 h_{\rm main}-\text{error}\ \le\ (V+2\delta)/d_0\ \le\ h_{\rm main}-m\,\mathrm{remainder}+\text{error}.
\]
The shared main height cancels. The loss from a missing fraction $4/b^2$ of the dimension and an index of size $O(d_0\log(1+b^2)\,\mathrm{budget})$ is absorbed. The source's choice of a second scale uses a long-window lemma; then division by $ma$ and the thick-place count gives the $x^{-1/2+8\epsilon}$ error (lines 149--172). The final step selects auxiliary quadratic twists with no zeros on a prescribed half-strip, so that a Dirichlet zero cannot be cancelled (lines 174--190 and \texttt{OneSidedCriterion.lean}).

\begin{center}\small
\begin{tabular}{p{0.29\linewidth}p{0.36\linewidth}p{0.26\linewidth}}
\toprule
Component & Inspected compact source (relative to \texttt{src/}) & Transfer status\\\midrule
Exact divisibility sum & \nolinkurl{New/Endgame/PrimeError.lean} & Finite difference proved below\\
Short and enough forms (A,B) & \nolinkurl{New/Endgame/ShortAndEnough.lean} & Needs new arithmetic weights\\
Long integral forms (C) & \nolinkurl{New/Endgame/LongForms.lean} & Needs same new lattice and index\\
Congruence control & \nolinkurl{New/Selmer/FitTwistsExist.lean}, \nolinkurl{New/Lift/ExponentTheorems.lean} & Specialized hypotheses, not plug-in results\\
Common Hecke traces & \nolinkurl{New/Trace/GoodFormTransfer.lean} & Potential tool; rigid labels and gap hypotheses\\
Mellin and real poles & \nolinkurl{Lib/Analytic/Landau/PolarLandau.lean} & Generalized in Section 3\\
Window selection & \nolinkurl{New/Endgame/OneSidedCriterion.lean} & Direct mathematical reuse for full prime powers\\
Quadratic twist avoidance & \nolinkurl{Lib/Analytic/QuadraticTwist/ZeroFreeTwist.lean} & Replaced by established abelian GRH and residue signs\\
\bottomrule
\end{tabular}\end{center}
The use of Selmer groups in C is substantial: the fitted twists must control all local denominators for the specific row representations and the first $m$ powers of a twist. The tame and crystalline length bounds in \texttt{ExponentTheorems} concern specific weights and local pairings. They do not directly bound the Tate--Shafarevich group of an arbitrary elliptic curve or arbitrary global Bloch--Kato Selmer groups. We therefore do not claim BSD, Goldbach, prime tuples, or general Artin holomorphy follows by invoking these names. The most direct viable transfer is to preserve the one-prime statistic and the main-height cancellation while changing the Frobenius representation it encodes.

\section{A one-sided finite-difference detector}
Let $c>1$ be fixed, and for $y>0$ define the nonnegative kernel
\[
 K_c(y)=\begin{cases}1-c^{-1},&0<y\le1,\\ y^{-1}-c^{-1},&1<y\le c,\\0,&y>c.\end{cases}
\]
Let $a(n)$ be real, supported on prime powers in the applications; assume $a(n)=O(\log(2n))$. Define
\begin{align}
 A(s)&=\sum_{n\ge2}a(n)n^{-s} \quad (\Rea s>1),\notag\\
 T(x)&=\sum_{n\le x}a(n)(1/n-1/x),\notag\\
 D_c(x)&=T(cx)-T(x)-\delta\log c
       =\frac1x\sum_{n\le cx}a(n)K_c(n/x)-\delta\log c.\label{defD}
\end{align}
The identity follows term by term, including at $n=x,cx$. It requires no value of a logarithmic-derivative constant at 1. We use the convention that positive residues are residues of $A(s)$ itself, not of $-A(s)$.

\begin{lemma}[Mellin identity]\label{mellin}
Suppose $A$ extends meromorphically to $\Rea s>\beta$ with $0<\beta<1$ and has at most a simple pole at 1, of residue $\delta$. For $\Rea s>1$,
\begin{align*}
 \int_1^\infty D_c(x)x^{-s}\,dx
 &=H_c(s)A(s)-\frac{\delta\log c}{s-1}-J_c(s),\\
 H_c(s)&=\frac{c^{s-1}-1}{s(s-1)},\qquad
 J_c(s)=c^{s-1}\int_1^c T(u)u^{-s}\,du.
\end{align*}
The term $J_c$ is entire. The pole at 1 is removable, $H_c(1)=\log c$, and $H_c$ has no zeros on $0<\Rea s<1$. On the positive real axis $H_c(\sigma)>0$ (including $\sigma=1$).
\end{lemma}
\begin{proof}
Absolute convergence gives
\[
 \int_1^\infty T(x)x^{-s}\,dx=\sum_n a(n)\int_n^\infty(1/n-1/x)x^{-s}\,dx=\frac{A(s)}{s(s-1)}.
\]
In the term with $T(cx)$, substitute $u=cx$; the omitted interval $[1,c]$ gives $J_c$. The remaining claims follow from the exponential function, since $c^{s-1}=1$ implies $\Rea s=1$, and $c^{\sigma-1}-1$ has the same sign as $\sigma-1$. This finite difference is important: the undifferenced kernel $1/[s(s-1)]$ is \emph{negative} on $(0,1)$ and would have the wrong residue sign for a deficit.
\end{proof}

\begin{theorem}[One-sided averaged detector]\label{landau}
In the setting of Lemma~\ref{mellin}, assume that all poles of $A$ in $\Rea s>\beta$ other than 1 are simple, lie in $\beta<\Rea s<1$, and have nonnegative residue whenever they are real. If for every real $\sigma>\beta$
\begin{equation}
 \int_1^\infty \pos{D_c(x)}x^{-\sigma}\,dx<\infty,\label{average}
\end{equation}
then $A$ has no pole in $\Rea s>\beta$ except possibly at 1. In particular it suffices that, for every $\epsilon>0$,
\begin{equation}
 \int_X^{2X}\pos{D_c(x)}\,dx=O_\epsilon(X^{\beta+\epsilon}).\label{dyadic}
\end{equation}
The pointwise condition $D_c(x)\le C_\epsilon x^{\beta-1+\epsilon}$ is stronger and also suffices.
\end{theorem}
\begin{proof}
Set $P=\pos{D_c}$ and $N=\pos{-D_c}$. The Mellin transform of $P$ is holomorphic on $\Rea s>\beta$, by dominated convergence on compact substrips (logarithmic derivatives are bounded by a slightly smaller power). The transform $F$ of $N$ converges for sufficiently large $\Rea s$, because the coefficient hypothesis gives $D_c(x)=O(\log(2x))$. By Lemma~\ref{mellin} and $D_c=P-N$, $F$ continues meromorphically to $\Rea s>\beta$. Each possible real pole there has nonpositive residue, and the point 1 is removable.

Let $\alpha$ be the real abscissa of convergence of the nonnegative transform of $N$. If $\alpha>\beta$, a simple pole at $\alpha$ would have strictly negative residue and force $F(\sigma)\to-\infty$ as real $\sigma\downarrow\alpha$, impossible. Hence $F$ is holomorphic at $\alpha$. Choose $\sigma_0>\alpha$ close enough that its Taylor disk reaches strictly to the left of $\alpha$. For $r>\sigma_0-\alpha$ within this disk, its Taylor series at $\sigma_0-r$ is
\[
 \sum_{k\ge0}\frac{r^k}{k!}\int_1^\infty N(x)(\log x)^kx^{-\sigma_0}\,dx.
\]
All summands are nonnegative. By monotone convergence its finite value is $\int_1^\infty N(x)x^{-(\sigma_0-r)}dx$, contradicting the definition of $\alpha$. (This is the usual Landau argument \cite{mv}.) Thus the transform of $N$ too is holomorphic in $\Rea s>\beta$. Lemma~\ref{mellin}, since $H_c$ is nonzero in the open critical strip, rules out all indicated poles, including nonreal ones. A dyadic decomposition proves that \eqref{dyadic} implies \eqref{average}, taking $\epsilon<\sigma-\beta$. Integrating the pointwise bound proves \eqref{dyadic}.
\end{proof}

\paragraph{Converse for the $L$-functions used here.} Assume their usual Euler products, meromorphic continuations and order-one completed functions with the standard gamma factors; assume all their zeros in the critical strip have real part $1/2$. The classical logarithmic derivative bound in any fixed strip $1/2+\epsilon\le\sigma\le2$ is $O_\epsilon(\log^2(3+|t|))$ for large $|t|$, with the pole at 1 excluded. One derivation is to subtract the Hadamard logarithmic derivative at $2+it$; near ordinates contribute $O_\epsilon(\log(3+|t|))$, using the usual zero count in unit intervals, and distant shells give a convergent sum with inverse-square decay. Gamma terms obey Stirling's bound. These are fixed-function estimates, not uniform estimates in the field or degree; see \cite{ik}. On the line 2 Mellin inversion of the reciprocal ramp gives $T(x)=(2\pi i)^{-1}\int_{(2)} A(s)x^{s-1}/[s(s-1)]\,ds$, by the elementary Perron kernel; hence $T(cx)-T(x)$ has multiplier $H_c(s)$. Shift to $1/2+\epsilon$ ($0<\epsilon<1/2$): $H_c=O((1+|t|)^{-2})$ on bounded vertical strips, and the only crossed pole contributes $\delta\log c$. The remaining integral is $O_\epsilon(x^{-1/2+\epsilon})$. This proves the reverse implication in our applications. For prime-only sums the omitted prime powers have $O(\log n)$ coefficients and their Euler series converges absolutely for $\sigma>1/2$. Equivalently their contribution to \eqref{defD} is $O_c(x^{-1/2}\log^2(2x))$: there are at most $O_c(\sqrt x\log x)$ powers $p^k\le cx$, $k\ge2$, each of weight $O(\log x)/x$. Finitely many bad primes (including their powers) contribute $O_c(\log^2(2x)/x)$. Thus replacing prime powers by primes changes neither implication at exponent $1/2$.

\section{Optimal nonnegative Frobenius deficits}\label{group}
Let $E/\Q$ be finite Galois, $G=\Gal(E/\Q)$, $H=[G,G]$, and $B=E^H$. The quotient $G/H$ is abelian, so $B$ is a subfield of a cyclotomic field. Dirichlet GRH applies to its zeta-function and to all characters of $G/H$. Let $F_j=E^{U_j}$ be subfields with no normality assumption on $U_j$. Let $\pi_j(g)=\#\mathrm{Fix}_{G/U_j}(g)$ be their permutation characters. Fix positive real numbers $b_j$ and set
\[
 \phi(g)=\sum_j b_j\pi_j(g),\qquad
 M(q)=\max_{gH=q}\phi(g),\qquad w(g)=M(gH)-\phi(g).
\]
This is a nonnegative real class function. It is the \emph{pointwise least} nonnegative function of the form $f(gH)-\phi(g)$: for each coset $q$, nonnegativity requires $f(q)\ge M(q)$. In particular it minimizes the sum over any collection of group elements with nonnegative weights, and its support is included in the support of every other such function. This optimization has no analytic assumption.

\begin{theorem}[Permutation-deficit criterion]\label{perm}
For unramified rational primes set $w_p=w(\Frob_p)$ and
\[
 \delta=\frac{1}{|G|}\sum_{g\in G}w(g),\qquad
 D(x)=\frac1x\sum_{p\le cx,\ p\text{ unramified in }E} w_p\log p\,K_c(p/x)-\delta\log c.
\]
Under Dirichlet GRH, the following are equivalent:
\begin{enumerate}[label=(\roman*)]
\item ERH for every $\zeta_{F_j}$, meaning all zeros in $0<\Rea s<1$ lie on $\Rea s=1/2$;
\item for every $\epsilon>0$, $D(x)=O_\epsilon(x^{-1/2+\epsilon})$;
\item for every $\epsilon>0$, $D(x)\le C_\epsilon x^{-1/2+\epsilon}$ for all sufficiently large $x$;
\item for every $\epsilon>0$, $\int_X^{2X}\pos{D(x)}dx=O_\epsilon(X^{1/2+\epsilon})$.
\end{enumerate}
\end{theorem}
\begin{proof}
Expand $M(q)=\sum_\chi m_\chi\chi(q)$ in characters of the finite abelian quotient. At unramified $p$, $a(p^k)=w(\Frob_p^k)\log p$ is bounded by a constant times $\log p$. The Artin permutation identity gives the Euler coefficients of $-\zeta'_{F_j}/\zeta_{F_j}$ as $\pi_j(\Frob_p^k)\log p$, since the cycle lengths are residue degrees. By Kronecker--Weber and the abelian Artin reciprocity identification, the full prime-power series equals, up to holomorphic finite Euler corrections on $\Rea s>0$,
\[
 A(s)=\sum_{\chi\in\widehat{G/H}}m_\chi\frac{-L'(s,\chi)}{L(s,\chi)}
      +\sum_j b_j\frac{\zeta'_{F_j}(s)}{\zeta_{F_j}(s)}.
\]
The corrections have no singularities on $\Rea s>0$ (the local eigenvalues have modulus at most one). The mean of each transitive permutation character is one. Thus the pole at 1 has residue $m_1-\sum_j b_j=\delta$. The known abelian RH means the first sum has no other pole right of $1/2$. At any zero $\rho$ of any $\zeta_{F_j}$ there, the residue of the second sum is $\sum_j b_j\operatorname{ord}_\rho\zeta_{F_j}>0$, with no cancellation. Classical Dedekind zeta-functions are nonzero on $\Rea s\ge1$ except for the pole at 1. Apply Theorem~\ref{landau} to the full prime-power series and use the negligible difference between powers and primes. This proves (iv) implies (i) by the functional equations of the zeta-functions. The converse argument after Theorem~\ref{landau} proves (i) implies (ii); the remaining implications are immediate. No general Artin holomorphy is needed.
\end{proof}

\begin{corollary}[Only nonsplit primes are needed]\label{split}
Take the regular action, normalized by $|G|$: $\phi(g)=\mathbf 1_{g=e}$. Then
\[
 w(g)=\mathbf 1_{g\in H\setminus\{e\}},\qquad
 \delta=\frac{1}{[B:\Q]}-\frac{1}{[E:\Q]}.
\]
Thus a single, unweighted, one-sided positive deficit test restricted to primes splitting completely in the maximal abelian subfield $B$ but not in $E$ is equivalent to ERH for $E$. Completely split primes themselves are absent from the test. The statement works equally with any specified abelian intermediate field by using its quotient in the construction (not necessarily the maximal one).
\end{corollary}
For a nontrivial abelian quotient the baseline is already known; in the abelian case the statistic vanishes identically. If $E/\Q$ is Galois and ERH for $E$ holds, all intermediate $\zeta_F$ satisfy ERH by the classical Aramata--Brauer entireness of $\zeta_E/\zeta_F$ (the extension $E/F$ is Galois). Brauer induction expresses each Artin $L$-function as a product with integer exponents of Hecke $L$-functions from cyclic subextensions in $E$. The latter are zero-free off the critical line inside the strip because the zeta of the corresponding abelian extension over its base is a product of those Hecke factors, each entire except a possible trivial pole at 1. Thus the detector also implies the absence of zeros \emph{and poles} of Artin $L$-functions off the critical line in the strip, but says nothing about possible poles \emph{on} the line. Artin holomorphy is a separate question.

\section{The Kummer fields of Artin's primitive-root conjecture}
Fix a nonzero integer $a$, not $-1$ or a square. For every squarefree $m\ge2$ let $B_m=\Q(\zeta_m)$ and $E_m=B_m(a^{1/m})$, with $h_m=[E_m:B_m]$ (a divisor of $m$). The extension $E_m/\Q$ is Galois as a splitting field of $X^m-a$. Corollary~\ref{split}, with the specified abelian subfield $B_m$, has the following direct arithmetic realization. Exclude the finite set of primes dividing $am$; for the rest,
\[
 w_{a,m}(p)=\begin{cases}
  1,&p\equiv1\pmod m\text{ and } a^{(p-1)/m}\not\equiv1\pmod p,\\
  0,&\text{otherwise},
 \end{cases}\qquad
 \delta_{a,m}=\frac{1-h_m^{-1}}{\varphi(m)}.
\]
Indeed splitting in $B_m$ is equivalent to the congruence. The cyclic group $\mathbb F_p^\times$ contains all $m$th roots of unity then; $X^m-a$ splits if and only if $a$ is an $m$th power, tested by the modular exponent. Frobenius fixes the roots exactly in that case. For an odd prime $\ell$ with $a$ not an $\ell$th power in $\Q$, $h_\ell=\ell$: an $\ell$th root in $B_\ell$ would, on taking norms, imply a rational $\ell$th power equal to $a^{\ell-1}$, forcing $a$ to be an $\ell$th power by valuations (and $\ell$ is odd). Over a field containing the roots of unity, the Kummer extension has degree $1$ or $\ell$ here. Hence $\delta_{a,\ell}=1/\ell$.

\begin{corollary}[Explicit reduction of Artin's conjecture]\label{artin}
If for every squarefree $m\ge2$ and every $\epsilon>0$ the quantity
\[
 \frac1x\sum_{p\le2x}w_{a,m}(p)\log p\,K_2(p/x)-\delta_{a,m}\log 2
 \ \le\ C_{a,m,\epsilon}x^{-1/2+\epsilon}
\]
for sufficiently large $x$, with the same conclusion also valid under the averaged condition of Theorem~\ref{perm}, then Artin's primitive-root conjecture holds for $a$, including the Hooley asymptotic density. The sum omits primes dividing $am$.
\end{corollary}
\begin{proof}
Each inequality implies ERH for $E_m$. Hooley's theorem \cite{hooley,moree} uses precisely GRH for the Dedekind zeta-functions of these squarefree cyclotomic--Kummer fields. Although our detector is fixed-field, the conclusion of each test is the exact zero statement. Hooley's uniform effective estimates then follow from the standard analytic theory under that zero statement, not from uniformity of the constants in our test. The case $m=1$ is ordinary RH, already known.
\end{proof}
This is a particularly direct analogue of the original nonsplit quadratic prime statistic: retain the congruence $p\equiv1\pmod m$, but replace the quadratic character by the $m$th power-residue test. It is \emph{not} a consequence of Dirichlet GRH alone, because the fields $E_m$ need not be abelian. For example, $a=2,m=3$ has $S_3$ splitting field, and the test consists exactly of primes $p\equiv1\pmod3$ for which 2 is not a cube. Proving the new inequalities even for all prime $m$ alone does not meet Hooley's squarefree-field hypothesis; the composite cases are required.

\section{A two-class detector for an \texorpdfstring{$S_4$}{S4} splitting field}\label{s4}
Consider $f(X)=X^4-X-1$, its splitting field $E$, and the natural four-point action of $S_4$. Its discriminant is $-283$; reduction modulo 2 is irreducible (no linear factor and no factor $X^2+X+1$), and modulo 7 it factors as $(X-3)(X^3+3X^2+2X+5)$, whose cubic has no root in $\mathbb F_7$. These patterns are also checked in the accompanying exact computations. Thus the Galois group contains a 4-cycle and a 3-cycle. Its order is divisible by 12; the only index-two subgroup is $A_4$, which excludes the odd 4-cycle, so the group is $S_4$. The maximal abelian subfield is the discriminant field $\Q(\sqrt{-283})$.

Let $F_4$ be the quartic field fixed by a point stabilizer and $F_6$ the field fixed by the order-four cyclic subgroup generated by a 4-cycle. The following character table entries are verified by the exact permutation calculation in \texttt{checks/check\_examples.py}. Here $\chi_3$ is the standard character, $\chi_{3'}=\chi_3\,\mathrm{sgn}$ and $\chi_2$ is the nontrivial two-dimensional character factoring through $S_4/V_4$.
\begin{center}\small
\begin{tabular}{lrrrrr}
\toprule
cycle type & $1^4$ & $1^2\,2$ & $2^2$ & $1\,3$ & $4$\\\midrule
class size & 1&6&3&8&6\\
$\chi_2$ & 2&0&2&$-1$&0\\
$\chi_3$ & 3&1&$-1$&0&$-1$\\
$\chi_{3'}$ & 3&$-1$&$-1$&0&1\\
$\pi_4$ &4&2&0&1&0\\
$\pi_6$ &6&0&2&0&2\\
$\phi=\pi_4+\pi_6$ &10&2&2&1&2\\
$M(\mathrm{sgn})$ &10&2&10&10&2\\
$w=M-\phi$ &0&0&8&9&0\\\bottomrule
\end{tabular}\end{center}
In particular $\delta=(3\cdot8+8\cdot9)/24=4$. For a prime $p\ne283$ the factorization type of $f$ modulo $p$ gives its Frobenius class. The functions $\pi_4=1+\chi_3$ and $\pi_6=1+\chi_2+\chi_{3'}$ can be checked by scalar products with the five irreducible characters.

All three nonlinear Artin factors of $S_4$ are entire: $\chi_2$ is induced from a nontrivial linear character of $A_4/V_4\simeq C_3$; the other two are induced from distinct linear characters of a stabilizer of a nontrivial character of $V_4$. More explicitly, the stabilizer has order 8 and is a split extension $V_4\rtimes C_2$; the invariant linear character of $V_4$ extends by either sign on the complement. The induced representations are irreducible by restriction to the three distinct $V_4$ characters permuted transitively, are of dimension 3, and differ by the sign twist. The nontrivial one-dimensional inducing characters give entire Hecke $L$-functions by induction and class field theory. The trivial and sign factors are already covered by RH and Dirichlet GRH. Thus the zeta-functions of $F_4,F_6$ together detect every remaining factor without cancellation.

\begin{theorem}[Sparse nonnegative $S_4$ test]\label{quartic}
ERH for $E$ is equivalent to the one-sided inequality (or averaged variant)
\[
 \frac1x\sum_{\substack{p\le2x\\p\ne283}}w_p\log p\,K_2(p/x)-4\log2
 \le C_\epsilon x^{-1/2+\epsilon}\qquad(\epsilon>0),
\]
where $w_p=8$ for two irreducible quadratic factors and $w_p=9$ for a linear and an irreducible cubic factor, and is zero otherwise. The unweighted version, $w_p=1$ on the same two types and $0$ elsewhere, also works, with centering $(11/24)\log2$ in place of $4\log2$.
\end{theorem}
\begin{proof}
The weighted version follows by Theorem~\ref{perm} for the two fields and the displayed character decompositions. Since both fields detect all nonlinear factors, the functional equations and product formula for $\zeta_E$ yield ERH. Conversely ERH for $E$ implies both subfield hypotheses by Aramata--Brauer. The unweighted version is Corollary~\ref{split}: the two types constitute $A_4\setminus\{e\}$, with density $11/24$.
\end{proof}

\begin{proposition}[Minimal support of strict-sign detectors]
Among nonnegative class functions on $S_4$ with strictly negative scalar products with all three nonlinear characters, the unique minimal support is the union of $2^2$ and $1\,3$. With only those classes in the support, writing the respective weights $u,v$, the exact condition is $u>0$ and $v>3u/4$.
\end{proposition}
\begin{proof}
The sum $\chi_3+\chi_{3'}$ equals $6,0,-2,0,0$ in the table order . Its strictly negative scalar product forces nonzero weight at $2^2$, even if weight at identity is allowed. The two-dimensional character is positive there and at identity, zero on the odd classes, and negative only at $1\,3$. Its strictly negative scalar product therefore forces nonzero weight at $1\,3$. On these two classes the inner products with the 3-dimensional characters are both $-3u/24$, and with $\chi_2$ the value is $(6u-8v)/24$. This proves both necessity and sufficiency. The minimality is for the specified strict-sign Fourier architecture, not a universal lower bound on all possible RH criteria.
\end{proof}
Compared with the signed tower detector in the earlier relative-races report, these weights do not need to observe completely split primes; the degree-four and degree-six realization gives a nonnegative target with all three unknown character directions and exact cancellation of the odd Frobenius strata. We do not assert that the underlying group identity or the analytic criterion has never appeared elsewhere.

\section{Beyond number fields: modular and symmetric-power RH}
The analytic theorem is not limited to finite Galois representations. Let $L(s,\pi)$ be a normalized, self-dual cuspidal automorphic $L$-function of degree $d$ over $\Q$, with real Euler traces $t(p^k)$ bounded by $d$ (tempered local parameters, omitting any finite exceptional set). Assume the usual continuation and functional equation, order-one completion and nonvanishing on $\Rea s\ge1$; these are standard for the holomorphic level-one examples below. Define $w(p^k)=d-t(p^k)\ge0$. Under ordinary RH,
\[
 A(s)=-d\frac{\zeta'(s)}{\zeta(s)}+\frac{L'(s,\pi)}{L(s,\pi)}
\]
up to finite harmless Euler corrections has residue $d$ at 1, and positive residue at each zero of $L(s,\pi)$ right of the critical line. Therefore the single inequality
\[
 \frac1x\sum_{p\le2x} (d-t(p))\log p\,K_2(p/x)-d\log2
 \le C_{\pi,\epsilon}x^{-1/2+\epsilon}
\]
for every $\epsilon>0$ is equivalent to RH for $L(s,\pi)$, also with the positive-part average. This follows exactly from Theorem~\ref{landau}, its converse, the prime-power tail, and the functional equation. For non-self-dual tempered functions one may instead use the product with the contragredient (real traces, no pole cancellation for entire cuspidal factors); we do not need that extension here.

A particularly explicit target is Ramanujan's discriminant form $\Delta(q)=q\prod_{n\ge1}(1-q^n)^{24}=\sum_{n\ge1}\tau(n)q^n$. In unitary normalization $t(p)=\tau(p)/p^{11/2}$, $d=2$ and Deligne's theorem \cite{deligne} bounds $|t(p)|\le2$. For each symmetric power $r\ge1$ of this non-CM level-one form, automorphy and cuspidality are established by Newton--Thorne \cite{nt}. Its normalized Satake parameters are $\alpha_p^{r-2j}$, $0\le j\le r$, where $|\alpha_p|=1$ and $\alpha_p+\alpha_p^{-1}=t(p)$. Thus $d=r+1$, $t_r(p)=\sum_{j=0}^r\alpha_p^{r-2j}$ is real and $|t_r(p)|\le r+1$. Equivalently $t_r(p)=U_r(t(p)/2)$, the Chebyshev polynomial of the second kind. For each fixed $r$, the detector with $w=r+1-t_r(p)$ is a precise nonnegative target for the symmetric-power RH, still beyond Dirichlet GRH. It does not prove Sato--Tate (already known), nor the Ramanujan bound for Maass forms (not supplied by Deligne), nor BSD's relation between central orders and ranks. In particular zero locations do not determine central orders.

\section{What would have to be proved geometrically}\label{geom}
\subsection{The exact divisibility increment}
For any integer $A\ge cx$ (not to be confused with the series $A(s)$), set
\[
 U_A(z)=\sum_{n\le A} a(n)(A/n-A/z)_+\quad(1\le z\le A).
\]
Then $U_A(z)=A T(z)$, and
\begin{equation}
 U_A(cx)-U_A(x)-A\delta\log c=A D_c(x).\label{increment}
\end{equation}
The term-by-term kernel is nonnegative for $a(n)\ge0$. Thus the proposed geometric target has exactly the finite positive sum shape of the original divisor counts, now in a difference at two nearby scales. Arithmetic progressions alone cannot supply the new weights: any fixed finite collection of Dirichlet characters depends on an abelian quotient of the absolute Galois group. A nonconstant function on $H$ (e.g. $w(e)=0,w(g)=1$ for $g\in H\setminus\{e\}$) cannot be represented this way, even after adjoining further abelian fields. To see the latter assertion, take the compositum with the abelian field; the commutator of its Galois group maps onto $[G,G]$ and is trivial on the abelian field. This is a strict information obstruction to a substitution using only fixed congruences. For Kummer fields, the residue-power test precisely carries the additional information.

\subsection{A two-scale certificate using the source window lemma}
For full prime-power coefficients (including finite Euler modifications), $D_c(e^v)$ has a globally meromorphic Mellin transform in $s$ by Lemma~\ref{mellin}; possible poles for $\Rea s>0$ are simple and none have real part at least 1. Growth $O(1+v)$ is more than sufficient for the source's \texttt{IsPolar} bound $O(e^v)$. The same holds for another such statistic $G_0(e^v)$, or for $G_0=0$. The source lemma \texttt{IsPolar.exists\_window\_sub\_mul\_le}, lines 73--100 of \texttt{New/Endgame/OneSidedCriterion.lean}, therefore supplies, for all sufficiently large $t$ and $0\le z\le e^t$, some $v\in[t,t^2]$ such that
\[
 D_c(e^v)-zG_0(e^v)\le e^{-t}.
\]
We use its mathematical content, not a locally replayed Lean instantiation. This step uses no extended RH.

\begin{proposition}[Sufficient transfer certificate]\label{cert}
Fix $L>1$. Suppose, for each $\epsilon>0$ and sufficiently large $x$, a coefficient $0\le z_x\le x^L$ can be fixed independently of $y$, and the geometric construction proves, for every $y$ in the window
\[
 x^L\le y\le\exp(L^2(\log x)^2),\qquad
 D_c(x)-D_c(y)+z_xG_0(y)\le C_\epsilon x^{-1/2+\epsilon}.
\]
Then the desired one-sided bound on $D_c(x)$ holds, and the corresponding RH in the preceding sections follows.
\end{proposition}
\begin{proof}
Apply the window lemma with $t=L\log x$ and $z=z_x$. For its chosen $y$, $D_c(y)-z_xG_0(y)\le x^{-L}$. Add the inequalities. The negligible $x^{-L}$ is absorbed in the stated bound.
\end{proof}
For a geometry-first strategy the input to this certificate should be a joint comparison of the \emph{same} covolume and integral index with the increment \eqref{increment} as its arithmetic remainder, and with a second-scale error. Two separate upper bounds at $x$ and $cx$ cannot be subtracted to obtain an upper bound on their difference. That tempting step is invalid. Either derive the increment directly from local divisibility or supply a matching lower bound for the subtracted term.

\subsection{Specific open obligations and milestones}
\begin{enumerate}[label=G\arabic*.]
\item \textbf{Local realization.} In the original proof the theta jets, formal groups of the chosen CM curves, and local divisibility counts encode $1-\chi_d$. Construct a replacement with the exact coefficients $w(\Frob_p^k)$ for a chosen field. Start with $a=2,m=3$, where the target has just one nonabelian character of degree two. A candidate finite-dimensional representation is the permutation action on the three roots of $X^3-2$; the arithmetic problem is to realize its normalized nonsplit increment in the relevant integral theta-lift volumes. Existence of that realization has not been shown here. For Kummer $m>3$, degree and conductor of the local model grow; for modular traces, nonintegral unitary normalization needs careful treatment rather than a literal exponent of a prime ideal.
\item \textbf{Short and enough.} Prove A and B for the same replacement family, with a positive row dimension, a missing proportion small enough to absorb the crude height, and an index estimate of the original order. The source's rank and level bookkeeping depends on its chosen data. A new proof needs existence and size estimates uniform throughout the second-scale window; importing a statement with the old \texttt{Data} type does not establish these.
\item \textbf{Long integral forms.} Prove the lower covolume bound with the same main height and index. The fitted-twist and local Selmer hypotheses must actually be checked for the new row representations; do not infer them from global Dirichlet GRH. This is a separate obstacle from G1 and likely the deepest after local realization.
\item \textbf{Quantified remainder and errors.} Establish exact identification with \eqref{increment} and the joint second-scale remainder, and prove that the total error after division by $m$ and the scale is $O(x^{-1/2+\epsilon})$ for each $\epsilon>0$. Check all signs, local ramified terms, and all quantifiers. The existing proof's $m\asymp\sqrt x$ is the origin of its saving, not an automatically portable asymptotic.
\end{enumerate}
An especially focused research milestone is G1 for the cubic Kummer field. It would distinguish the transpositions, 3-cycles and identity while the known abelian sign character only distinguishes even from odd. The unweighted detector in Corollary~\ref{split} then requires only the even nonidentity class. If G1--G4 succeed for every fixed squarefree $m$, Corollary~\ref{artin} closes Artin's conjecture via Hooley. For the modular problem, G1 must encode algebraic Hecke traces with a normalization; standard results on their size do not give the required divisor identity. No such geometric estimate is proved in this memorandum.

\section{Reproduction, assessment and limits}
The accompanying Python checks were run. They verify the $S_4$ coset character, inner products and minimal-support weights exactly; the affine groups for $\ell=3,5,7,11,13$; the rational kernel identity; and finite-field factorization patterns for $X^4-X-1$ at all primes up to 2000 except 283. The CSV diagnostics for cubic and quintic Kummer fields and the quartic detector use ordinary floating-point logarithms and are illustrative only. No finite computation can establish the asymptotic inequality without a separate analytic bound.

\begin{center}\small
\begin{tabular}{p{0.33\linewidth}p{0.55\linewidth}}
\toprule
Target & Outcome and remaining risk\\\midrule
Nonabelian ERH & Proved optimal nonnegative permutation-deficit detector; G1--G4 open.\\
Artin primitive roots & Explicit modular exponent criterion per squarefree Kummer field; Hooley reduction rigorous; estimates open.\\
Modular / symmetric-power RH & Proved positive trace-deficit detector under known analytic input; arithmetic realization open.\\
General Artin holomorphy & Off-line zero/pole exclusion would follow from ERH of Galois field; central-line holomorphy not addressed.\\
BSD and general Bloch--Kato & Congruence-control ideas are plausible local tools, but no general Selmer estimate follows.\\
Prime-pair conjectures & Source controls one prime-power statistic; pair correlations are not implied.\\\bottomrule
\end{tabular}\end{center}
The analytic and finite-group assertions are proved here in conventional mathematics using cited standard facts. The source-level map is a selected-file reading, not a dependency-closure audit. The new proofs are not yet formalized in Lean; the supplied source-audit script is for later local use on a checkout, and was not run against the upstream checkout here. The GitHub connector did supply the file contents and blob IDs, and the earlier archives supplied selected text. No repository was modified. No new major conjecture is claimed resolved. The novelty of the combination, particularly the finite-difference sign repair and the optimization, is a research proposal supported by the proofs but not a certified literature-priority claim.

\begin{thebibliography}{99}\small
\bibitem{repo} L. Sernau et al., \emph{GRH for Dirichlet L-functions, in Lean 4}, 25 Nov. 2026. \url{https://github.com/lsernau/grh-lean}. Compact \texttt{README}, \texttt{GUIDE}, \texttt{Statements}, \texttt{New/Main} and files identified in the source audit.
\bibitem{cert} Lean FRO and Mathlib Initiative, \emph{Lean verification of the Riemann Hypothesis}, 21 Nov. 2026. \url{https://leanprover-community.github.io/blog/posts/riemann-hypothesis/}.
\bibitem{rubinstein} M. Rubinstein, \emph{Alp\"oge/Claude proof of the Generalized Riemann Hypothesis}, verification and working exposition, including shortened edition status. \url{https://www.math.uwaterloo.ca/~mrubinst/alpoge_claude_grh/}.
\bibitem{prior} ProveIt incoming research packets, directory accessed 25 Nov. 2026. \url{https://github.com/VladimirReshetnikov/ProveIt/tree/main/docs/incoming}. Archive names and Git blob IDs in \texttt{prior-audit.tsv}.
\bibitem{mv} H. L. Montgomery and R. C. Vaughan, \emph{Multiplicative Number Theory I: Classical Theory}, Cambridge, 2007, especially Ch. 15 (Landau and oscillation theorems).
\bibitem{neukirch} J. Neukirch, \emph{Algebraic Number Theory}, Springer, 1999. Artin reciprocity, Dedekind zeta, and Artin $L$-functions.
\bibitem{brauer} R. Brauer, On the zeta-functions of algebraic number fields, \emph{Amer. J. Math.} 69 (1947), 243--250. Aramata--Brauer entireness theorem; Brauer induction and its use for Artin $L$-functions are also classical.
\bibitem{hooley} C. Hooley, On Artin's conjecture, \emph{J. Reine Angew. Math.} 225 (1967), 209--220. \url{https://doi.org/10.1515/crll.1967.225.209}.
\bibitem{moree} P. Moree, Artin's primitive root conjecture---a survey, \emph{Integers} 12 (2012), A13. \url{https://arxiv.org/abs/math/0412262}.
\bibitem{ik} H. Iwaniec and E. Kowalski, \emph{Analytic Number Theory}, AMS Colloquium Publications 53, 2004. Classical zero counts, logarithmic derivatives and automorphic $L$-functions.
\bibitem{deligne} P. Deligne, La conjecture de Weil I, \emph{Publ. Math. IH\'ES} 43 (1974), 273--307.
\bibitem{nt} J. Newton and J. A. Thorne, Symmetric power functoriality for holomorphic modular forms, \emph{Publ. Math. IH\'ES} 134 (2021), 1--116; and II, 117--152. \url{https://doi.org/10.1007/s10240-021-00127-3}, \url{https://doi.org/10.1007/s10240-021-00126-4}.
\bibitem{suzuki} M. Suzuki, On variants of Chebyshev's conjecture, 2024. \url{https://arxiv.org/abs/2411.07436}. Related one-sided criteria and prime races.
\end{thebibliography}
\end{document}
